\documentclass[11pt,a4paper]{article}
\usepackage[margin=2.4cm]{geometry}
\usepackage{fontspec}
\usepackage{xcolor}
\usepackage{titlesec}
\usepackage{enumitem}
\usepackage{booktabs}
\usepackage{tabularx}
\usepackage{longtable}
\usepackage{graphicx}
\usepackage{fancyhdr}
\usepackage{hyperref}
\usepackage{listings}
\usepackage{parskip}

\definecolor{ink}{HTML}{0F172A}
\definecolor{accent}{HTML}{1D4ED8}
\definecolor{soft}{HTML}{475569}
\definecolor{panel}{HTML}{F1F5F9}

\setmainfont{Times New Roman}
\setsansfont{Segoe UI}
\setmonofont{Consolas}

\hypersetup{colorlinks=true,linkcolor=accent,urlcolor=accent,bookmarksnumbered=true}
\titleformat{\section}{\sffamily\Large\bfseries\color{ink}}{\thesection}{0.6em}{}
\titleformat{\subsection}{\sffamily\large\bfseries\color{ink}}{\thesubsection}{0.6em}{}
\titleformat{\subsubsection}{\sffamily\normalsize\bfseries\color{soft}}{\thesubsubsection}{0.6em}{}
\titlespacing*{\section}{0pt}{1.6em}{0.6em}

\pagestyle{fancy}
\fancyhf{}
\renewcommand{\headrulewidth}{0.4pt}
\fancyhead[L]{\footnotesize\sffamily\color{soft}JC2001 Software Engineering \textperiodcentered{} Group 5}
\fancyhead[R]{\footnotesize\sffamily\color{soft}StudyGraph User Manual}
\fancyfoot[C]{\footnotesize\sffamily\color{soft}\thepage}

\lstset{
  basicstyle=\ttfamily\footnotesize,
  backgroundcolor=\color{panel},
  frame=none,
  breaklines=true,
  columns=fullflexible,
  xleftmargin=0.4em,
  xrightmargin=0.4em,
  aboveskip=0.6em,
  belowskip=0.8em,
  showstringspaces=false
}

\newcommand{\code}[1]{\texttt{\small\detokenize{#1}}}

\begin{document}

\begin{titlepage}
\centering
\vspace*{3.2cm}
{\sffamily\small\color{soft} JC2001 Software Engineering \textperiodcentered{} Group 5}\\[0.5cm]
{\sffamily\bfseries\fontsize{34}{40}\selectfont\color{ink} StudyGraph}\\[0.35cm]
{\sffamily\large\color{soft} Graph-Augmented Misconception Diagnosis\\ for Exam Preparation}\\[2.2cm]
{\sffamily\bfseries\Large\color{accent} User Manual}\\[0.3cm]
{\sffamily\normalsize\color{soft} Installation, operation, and interface reference for the Cycle 3 proof of concept}\\[3.2cm]

\begin{minipage}{0.78\textwidth}
\centering
{\sffamily\small\color{soft}
Document version 1.0 \textperiodcentered{} 26 November 2026\\[0.2cm]
Prepared for examiners, demonstrators, and future maintainers\\
University of Aberdeen \textperiodcentered{} South China Normal University}
\end{minipage}

\vfill
{\sffamily\footnotesize\color{soft} This manual describes the software as submitted at the Cycle 3 proof-of-concept HEAD.\\
It is self-contained: no external account, API key, or network connection is required to follow it.}
\end{titlepage}

\tableofcontents
\newpage

\section{About This Manual}

StudyGraph is a local, single-process prototype that helps university students revise rule-dense courses.
A student answers a multiple-choice question; the system identifies the \emph{specific misconception} behind a
wrong answer, shows the missing prerequisite, and asks a Socratic question instead of revealing the answer.

This manual is written for three readers:

\begin{itemize}[leftmargin=1.4em,itemsep=0.15em]
  \item \textbf{Examiners and demonstrators} who want to run the prototype and try it in a few minutes. Start with Section~\ref{sec:start}.
  \item \textbf{The development team} who needs a reference for the five console views and the REST contract. See Sections~\ref{sec:tour}--\ref{sec:api}.
  \item \textbf{Future maintainers} who need to know what is deliberately out of scope. See Section~\ref{sec:limits}.
\end{itemize}

The manual describes the software \emph{as delivered}: it does not promise features that exist only on the roadmap.
Where the implementation is narrower than the design, the manual says so explicitly.

\subsection{What the prototype does}

\begin{tabularx}{\textwidth}{@{}lX@{}}
\toprule
\textbf{Capability} & \textbf{Behaviour in this build} \\
\midrule
Curated question bank & 17 questions across 9 subjects, each with exactly four options and three distractor traps (51 traps total) \\
Trap diagnosis & A wrong answer returns the trap it selected: title, description, prerequisite, and radar dimension \\
Local concept graph & The four-field model (\code{GraphData}, \code{GraphNode}, \code{GraphEdge}, \code{TrapDetail}) is returned with the diagnosis \\
Socratic tutor & A short probing question assembled from reviewed evidence; no answer text is disclosed \\
Learning analytics & Five-dimension mastery radar plus a session trend line, computed in the browser \\
Deterministic offline mode & All of the above work with no API key and no network connection \\
\bottomrule
\end{tabularx}

\subsection{What the prototype deliberately is not}

\begin{itemize}[leftmargin=1.4em,itemsep=0.15em]
  \item Not a multi-user service. Authentication is a demonstration login; there is no server-side session enforcement.
  \item Not a production answer-key service. Question payloads expose the correct option; the design contract is documented as a known limitation.
  \item Not an online-generation product at runtime. The large-language-model path is optional and is used only to rephrase an already-decided explanation.
  \item Not a mobile application. The console is a responsive browser interface served by the same process.
\end{itemize}

\section{System Requirements}
\label{sec:req}

\begin{tabularx}{\textwidth}{@{}lX@{}}
\toprule
\textbf{Item} & \textbf{Requirement} \\
\midrule
Operating system & Windows 10/11, macOS, or Linux; verified on Windows 11 \\
Python & 3.10 or later (verified on 3.13/3.14); \code{venv} recommended \\
Packages & The twelve dependencies in \code{requirements.txt}; a network connection is needed once, at install time \\
Browser & Any current browser with ES6 and CSS custom properties (Chrome, Edge, Firefox, Safari) \\
Port & \code{8000} by default; overridable with \code{--port} or \code{PORT} \\
Disk & About 3~MB for the packaged source, plus the environment \\
Network & Not required at run time; required only for the first \code{pip install} and for the optional model probe \\
\bottomrule
\end{tabularx}

\section{Getting Started}
\label{sec:start}

The procedure is five commands and takes about a minute once the dependencies are cached.
Run everything from the project root (the directory containing \code{backend/} and \code{frontend/}).

\subsection{Install and start}

\begin{lstlisting}[language=bash]
# 1. create and activate an isolated environment
python -m venv .venv
.venv\Scripts\activate          # Windows
source .venv/bin/activate       # macOS / Linux

# 2. install the dependencies declared in the repository
pip install -r requirements.txt

# 3. start the service (the project root, not the backend folder)
uvicorn backend.app.main:app --port 8000

# 4. open the console
start http://127.0.0.1:8000     # or paste the address into a browser
\end{lstlisting}

The service prints the bound address, for example
\code{Uvicorn running on http://127.0.0.1:8000}. Opening that address serves the study console;
\code{http://127.0.0.1:8000/docs} serves the interactive OpenAPI documentation.

\begin{quote}
\small\color{soft}\textbf{Common first mistake.} Run the server from the project root. Running
\code{python backend/run.py} from that same root fails with
\code{ModuleNotFoundError: No module named 'backend'}, because Python places the script's own directory on
\code{sys.path}. Use the \code{uvicorn backend.app.main:app} command above instead.
\end{quote}

\subsection{Verify the installation}

Three checks confirm that the system is healthy before a demonstration:

\begin{lstlisting}[language=bash]
curl http://127.0.0.1:8000/api/health
curl http://127.0.0.1:8000/api/quiz?limit=1
python -m pytest tests/ -q
\end{lstlisting}

The first returns a status payload with \code{"status":"ok"}; the second returns a question;
the third runs the offline regression suite and should end with
\code{all tests passed} (no API key and no network required).

\subsection{Running completely offline}

Create an \code{.env} file only if you want optional integrations; the server starts without one.
When \code{DASHSCOPE_API_KEY} (or \code{QWEN_API_KEY}) is unset, the interface reports the model path as
\emph{not configured} and every diagnosis is served from the curated local bank. This is the default and
preferred configuration for marking, because it makes results reproducible.

\section{Interface Tour}
\label{sec:tour}

The console has five views, listed in the sidebar. Each view is a single page section; no build step,
package manager, or bundler is involved.

\begin{center}
\small
\begin{tabularx}{\textwidth}{@{}lX@{}}
\toprule
\textbf{View} & \textbf{Purpose} \\
\midrule
\code{view-home} & Landing summary: bank size, subjects, and the entry points to practice and tutor \\
\code{view-quiz} & Diagnostic practice: question, four options, hint, submit, and the trap card \\
\code{view-tutor} & Socratic dialogue: follow-up questions grounded in the same reviewed evidence \\
\code{view-graph} & Concept graph: the local subgraph for the current question \\
\code{view-report} & Learning analytics: five-dimension radar and session trend \\
\bottomrule
\end{tabularx}
\end{center}

A theme control in the header switches between light and dark presentation; the choice is stored in the browser
and does not affect the data.

\subsection{Reading a question}

Each question card shows the subject, a review-state badge, a source reference, the stem, and four options.
Two controls matter:

\begin{itemize}[leftmargin=1.4em,itemsep=0.15em]
  \item \textbf{View Hint} reveals a one-line nudge that narrows the choice without naming the answer.
  \item \textbf{Submit Answer} records the attempt and opens the feedback panel.
\end{itemize}

\subsection{Reading the trap card}

After a wrong answer the feedback panel shows four fields, which are exactly the fields the diagnosis endpoint returns:

\begin{center}
\small
\begin{tabularx}{\textwidth}{@{}lX@{}}
\toprule
\textbf{Field} & \textbf{Meaning} \\
\midrule
Trap title & The named misconception the selected distractor encodes \\
Description & Why that misconception is tempting, in one paragraph \\
Prerequisite & The concept the student is missing and should revise first \\
Radar dimension & Which of the five mastery dimensions this attempt moves \\
\bottomrule
\end{tabularx}
\end{center}

The correct option is highlighted at the same time. The trap card is the core interaction of the prototype:
it converts a binary right/wrong mark into a named, actionable gap.

\subsection{Using the Socratic tutor}

The tutor view accepts a free-text follow-up. Replies are assembled from the reviewed evidence attached to the
current question, so they stay inside the curriculum. If an external model is configured, it is asked only to
rephrase that evidence more conversationally; the content does not change.
If no model is configured, or the call fails or times out, the tutor falls back to the same reviewed evidence and
labels the reply as local.

\subsection{Reading the analytics}

The report view renders two visualisations from attempt history stored in the browser:

\begin{itemize}[leftmargin=1.4em,itemsep=0.15em]
  \item A five-axis radar over \emph{recall}, \emph{boundary}, \emph{trap defence}, \emph{synthesis}, and \emph{precision}.
  \item A session trend line comparing accuracy across attempts.
\end{itemize}

With no attempts recorded the radar is empty and the trend is flat; this is expected on a fresh browser profile.

\section{Worked Example: One Diagnosis End to End}

This section follows a single attempt through the system, which is the fastest way to understand the architecture.

\begin{enumerate}[leftmargin=1.6em,itemsep=0.3em]
  \item Open the practice view and choose a question, for example a question on statutory periods in civil law.
  \item Select a distractor, for example the option that applies the ordinary limitation period to a formative right.
  \item Press \textbf{Submit Answer}. The console posts to \code{POST /api/diagnose} with the question identifier and the selected key.
  \item The router validates the body with a Pydantic model, then delegates to the diagnostic service.
  \item The service looks up the option, reads its correctness flag, and --- on a wrong option --- reads the attached trap record.
  \item The response carries the trap title, description, prerequisite, radar dimension, the local concept subgraph, and a Socratic prompt.
  \item The console renders the trap card, highlights the correct option, appends the attempt to browser storage, and redraws the radar.
\end{enumerate}

There is no model inference on this path. The diagnosis is a dictionary lookup followed by field assembly, which is
why it is deterministic and why it works offline.

\section{API Reference}
\label{sec:api}

All endpoints are mounted under \code{/api}. Requests and responses are JSON; responses use camelCase field names.
The interactive documentation at \code{/docs} lists the exact schemas.

\subsection{Endpoint summary}

\begin{center}
\small
\begin{tabularx}{\textwidth}{@{}llX@{}}
\toprule
\textbf{Method} & \textbf{Path} & \textbf{Purpose} \\
\midrule
GET & \code{/api/health} & Liveness and configuration summary \\
GET & \code{/api/health/llm} & Optional model probe; masks credentials \\
POST & \code{/api/quiz} & Question list, optionally filtered by subject \\
POST & \code{/api/diagnose} & Trap diagnosis for one selected option \\
GET & \code{/api/reports/summary} & Five-dimension mastery summary \\
GET & \code{/api/reports/trend} & Attempt trend for a subject \\
GET & \code{/api/auth/demo-users} & Demonstration personas \\
POST & \code{/api/auth/login} & Demonstration login \\
\bottomrule
\end{tabularx}
\end{center}

\subsection{Example requests}

\begin{lstlisting}[language=bash]
# health
curl http://127.0.0.1:8000/api/health

# questions for one subject
curl -X POST http://127.0.0.1:8000/api/quiz \
     -H "Content-Type: application/json" \
     -d '{"subject":"law","limit":5}'

# diagnose a selected option
curl -X POST http://127.0.0.1:8000/api/diagnose \
     -H "Content-Type: application/json" \
     -d '{"question_id":"law_1","selected_option":"B"}'

# mastery summary and trend
curl "http://127.0.0.1:8000/api/reports/summary?subject=law"
curl "http://127.0.0.1:8000/api/reports/trend?subject=all"
\end{lstlisting}

\subsection{Error behaviour}

Invalid identifiers return \code{404}; structurally invalid bodies and unprocessable selections return \code{422}
with a field-level diagnostic. The suite in \code{tests/test_adversarial_backend.py} proves that malformed input
never produces a \code{500}: injection strings, type confusion, overflow payloads, and concurrent bursts all return
a controlled status.

\section{Troubleshooting}

\begin{tabularx}{\textwidth}{@{}lX@{}}
\toprule
\textbf{Symptom} & \textbf{Cause and remedy} \\
\midrule
\code{ModuleNotFoundError: backend} & The server was started from the wrong directory. Run \code{uvicorn backend.app.main:app} from the project root. \\
Port already in use & Another process holds \code{8000}. Start with \code{--port 8001} and open the matching address. \\
Blank charts in the report view & No attempts recorded yet, or the browser blocked local storage. Answer one question first. \\
\code{/api/health/llm} reports not configured & No model key is present. This is the intended offline configuration. \\
\code{401} on a report call & The request omitted the demonstration identity header the console normally sends. \\
Slow first start & The first import compiles the application; later starts are effectively instant. \\
\bottomrule
\end{tabularx}

\section{Limitations and Scope}
\label{sec:limits}

The following limits are deliberate and are stated here so that no reader over-estimates the prototype:

\begin{itemize}[leftmargin=1.4em,itemsep=0.2em]
  \item \textbf{One user at a time.} Attempt history lives in browser storage, not in a shared database.
  \item \textbf{Demonstration authentication.} Passwords are salted and hashed, but there is no server-side session enforcement.
  \item \textbf{Answer keys in question payloads.} The quiz response includes correctness flags; a production build would withhold them.
  \item \textbf{Wildcard CORS.} Acceptable for a local demonstration, not for a public deployment.
  \item \textbf{Offline-first evaluation.} The GraphRAG benchmark runs as a separate offline pipeline; the runtime does not query the graph database.
\end{itemize}

\section{Where to Look Next}

\begin{tabularx}{\textwidth}{@{}lX@{}}
\toprule
\textbf{Question} & \textbf{Document} \\
\midrule
How is the system designed? & \code{reports/final_report/StudyGraph\_Cycle3\_Final\_Report.pdf}, Chapters 3--5 \\
How is it tested? & The same report, Chapter 6, and Appendix C \\
How do I run the offline evaluation? & \code{graphrag\_final/README.md} \\
What are the exact HTTP contracts? & \code{http://127.0.0.1:8000/docs} at run time \\
How do I contribute? & \code{README.md} at the repository root \\
\bottomrule
\end{tabularx}

\vspace{1.2em}
\begin{quote}
\small\color{soft}
If a step in this manual does not reproduce, treat that as a defect in the manual and report it rather than
working around it silently: the manual is part of the deliverable and is expected to match the submitted code.
\end{quote}

\end{document}
